\section*{第 \ref{ch:b2:coordination-complexes} 章——过渡金属与配合物}

\begin{solution}{exo:b2:coordination-complexes:1}
\ce{Ti^3+} $d^1$, \ce{V^3+} $d^2$, \ce{Cr^3+} $d^3$, \ce{Mn^2+} $d^5$, \ce{Fe^3+} $d^5$, \ce{Co^2+}
$d^7$, \ce{Ni^2+} $d^8$, \ce{Zn^2+} $d^{10}$.
\end{solution}

\begin{solution}{exo:b2:coordination-complexes:2}
水 1，en 2，乙二酸根 2，EDTA 6，硫氰酸根 1；硫氰酸根是
两可配体（通过 S 或 N 配位）。
\end{solution}

\begin{solution}{exo:b2:coordination-complexes:3}
硫酸四氨合铜(II)；六氰合铁(III)酸钾；
氯化二氯四氨合钴(III)；四羰基合镍(0)。
\end{solution}

\begin{solution}{exo:b2:coordination-complexes:4}
直线形；四面体形；平面正方形（$d^8$，第三过渡系）；八面体形。
\end{solution}

\begin{solution}{exo:b2:coordination-complexes:5}
两种：顺式（两个氯相邻）和反式（两个氯相对）。四面体中
只会有一种。
\end{solution}

\begin{solution}{exo:b2:coordination-complexes:6}
\ce{Cr(CO)6} 18; \ce{Fe(CO)5} 18; \ce{Ni(CO)4} 18; \ce{V(CO)6}: $5 + 12 = 17$;
\ce{Mn2(CO)10} 每个锰 18 个电子。\ce{V(CO)6} 违背该规则：它具有顺磁性，并且
很容易被还原为 \ce{[V(CO)6]-}，后者的电子数为 18。
\end{solution}

\begin{solution}{exo:b2:coordination-complexes:7}
二茂铁：Fe(II) $d^6$ + 12 = 18。二茂钴：Co(II) $d^7$ + 12 = 19。二茂钴易
失去多余的电子，生成二茂钴阳离子 \ce{[Co(C5H5)2]+}，其电子数为 18。
\end{solution}

\begin{solution}{exo:b2:coordination-complexes:8}
一个 EDTA 取代六个水分子：\ce{[Ni(H2O)6]^2+ + edta^4- -> [Ni(edta)]^2- +
6H2O}，由两个粒子变为七个粒子。反应熵变很大且为正，
$K$ 非常大。六个独立配体取代六个水分子时，
粒子数没有增加。
\end{solution}

\begin{solution}{exo:b2:coordination-complexes:9}
亚硝酸根-$\kappa N$（\ce{Co-NO2}，硝基异构体，黄色）和亚硝酸根-$\kappa O$
（\ce{Co-O-N=O}，亚硝酸根异构体，红色）。
\end{solution}

\begin{solution}{exo:b2:coordination-complexes:10}
三种：反式异构体（非手性，有镜面）和顺式异构体的
两个对映异构体（两个螯合物和两个氯以一种或另一种方式缠绕）。
\end{solution}

\begin{solution}{exo:b2:coordination-complexes:11}
平面六边形：B 位于 1,2（邻位）、1,3（间位）、1,4（对位），共三种异构体。
三棱柱：B 沿三角形的一条边、沿竖直棱，或跨越一个
矩形面：三种。八面体给出两种。尽管
多次尝试、研究多种化合物，仍只找到两种，由此排除另外两种几何构型——前提是
没有遗漏第三种异构体。
\end{solution}

\begin{solution}{exo:b2:coordination-complexes:12}
\ce{[RhCl(PPh3)3]}: Rh(I) $d^8$ + 4 × 2 = 16. \ce{[IrH(CO)(PPh3)3]}: Ir(I) $d^8$ + 5 × 2 =
18。铑配合物有 16 个电子，还可再接受两个电子（氧化加成，
见\cref{ch:b2:catalytic-cycles}）。
\end{solution}

\begin{solution}{pb:b2:coordination-complexes:1}
\textbf{1.} 配位层外的氯离子，在溶液中自由存在。
\textbf{2.} A：0；B：1；C 和 D：2。
\textbf{3.} A：$[\ce{Co(NH3)6}]^{3+}$ + 3 \ce{Cl-}：4 个离子；B：3；C、D：2。结果一致。
\textbf{4.} +3。
\textbf{5.} $d^6$。
\textbf{6.} 每种均为六配位：6 \ce{NH3}；5 \ce{NH3} + 1 Cl；4 \ce{NH3} + 2 Cl。
\textbf{7.} \ce{[Co(NH3)6]Cl3}；\ce{[CoCl(NH3)5]Cl2}；\ce{[CoCl2(NH3)4]Cl}（出现两次）。
\textbf{8.} 氯化六氨合钴(III)；氯化一氯五氨合钴(III)。
\textbf{9.} 二氯四氨合钴(III)。
\textbf{10.} 几何（顺–反）异构体。
\textbf{11.} \ce{[CoCl3(NH3)3]}，中性配合物：不含离子，
因而不会立即沉淀出氯离子。
\textbf{12.} 两种：面式和经式。
\textbf{13.} 顺式：两个 Cl 的夹角为 $90^\circ$；反式：为 $180^\circ$。
\textbf{14.} 三种（六边形上的邻、间、对位）。
\textbf{15.} 三种。
\textbf{16.} 它支持八面体构型，是三者中唯一预测为两种的。
\textbf{17.} 否：没有第三种异构体只是否定性的证据。若找到第三种异构体，
则会排除八面体构型。
\textbf{18.} 形成螯合物的那一种：双齿乙二酸根只能跨越
两个顺式位置。
\textbf{19.} 钴周围有三个 en 螯合物，每个都跨越两个顺式位置。
\textbf{20.} 两者都不是：既无镜面，也无反演中心。
\textbf{21.} 两种对映异构体，$\Delta$ 和 $\Lambda$。
\textbf{22.} 六个配体按三维方式排列，如
八面体那样；平面排列则会得到非手性离子。
\textbf{23.} 顺式：手性（两种对映异构体）。反式：非手性。
\textbf{24.} 它表明旋光性一般属于分子几何构型，
而非仅与碳有关，并证实了八面体构型。
\textbf{25.} 八面体预测 \ce{[CoCl2(NH3)4]+} 有 $\boldsymbol{2}$ 种异构体；
平面六边形和三棱柱则预测有 $\boldsymbol{3}$ 种。
\end{solution}
